\chapter{Introducción}

    \section{Antecedentes}

        La descarga electrostática (ESD, por sus siglas en inglés) representa una de 
        las principales causas de daño en dispositivos electrónicos sensibles, 
        especialmente en la industria de semiconductores. A medida que los circuitos integrados 
        han reducido su tamaño, su susceptibilidad a eventos ESD ha incrementado, lo que ha 
        llevado a la implementación de estrictos controles en ambientes de manufactura y laboratorios electrónicos \cite{versalogic_esd}.

        Para mitigar estos efectos, se han desarrollado técnicas de control basadas en la 
        puesta a tierra del personal, utilizando dispositivos como pulseras antiestáticas y calzado conductivo. 
        Estos elementos permiten disipar las cargas acumuladas en el cuerpo humano, reduciendo el riesgo de descargas 
        hacia componentes electrónicos \cite{STM97_ANSI}.

        Los sistemas de verificación ESD han evolucionado desde dispositivos simples que únicamente 
        indican condiciones de "pasa" o "falla", hasta sistemas más avanzados que integran múltiples funcionalidades. 
        Los testers básicos de pulseras antiestáticas, como el Desco 19240, permiten validar la continuidad eléctrica 
        del sistema de puesta a tierra, pero no ofrecen información detallada ni capacidades de integración con otros sistemas.

        Por otro lado, soluciones más complejas como el SmartLog Pro 2 o sistemas de control ESD con acceso integrado 
        incluyen características como identificación de usuario mediante RFID, registro de pruebas, comunicación en red 
        y generación de reportes. Estos sistemas permiten mejorar la trazabilidad y el cumplimiento de normas industriales, 
        pero presentan limitaciones importantes como su alto costo y baja flexibilidad de personalización \cite{SmartLogPro2}.

        En la literatura técnica también se han propuesto diferentes enfoques para el control ESD, incluyendo sistemas 
        de monitoreo continuo mediante IoT y medición de campos electrostáticos en líneas de producción. Sin embargo, 
        estos trabajos no abordan de forma integral el diseño de sistemas accesibles, personalizables y orientados a 
        aplicaciones de laboratorio.

        Adicionalmente, desde el punto de vista de instrumentación electrónica, la medición de la resistencia cuerpo–tierra 
        implica desafíos asociados a la medición de altas impedancias, donde factores como corrientes de fuga, ruido eléctrico 
        y condiciones ambientales pueden afectar significativamente la precisión \cite{keithley_lowlevel}.

        Este contexto evidencia la necesidad de desarrollar soluciones que integren medición precisa, 
        interacción con el usuario, monitoreo ambiental y capacidades de registro, manteniendo un bajo costo y 
        alta adaptabilidad a distintos entornos.

    \section{Planteamiento del Problema}
        En los laboratorios de electrónica y manufactura de semiconductores, 
        el control ESD se realiza mediante sistemas de verificación de resistencia cuerpo a tierra. 
        Sin embargo, del análisis de las soluciones comerciales disponibles, no se identifican sistemas de bajo costo 
        que integren de manera conjunta medición cuantitativa, identificación de usuario, registro de datos 
        y capacidad de personalización.

        Las soluciones de bajo costo se limitan a verificaciones de tipo pasa/falla, mientras que los sistemas 
        avanzados que incorporan estas funcionalidades presentan un costo elevado y arquitectura cerrada, 
        lo que restringe su adaptación a necesidades específicas.

        Esta situación evidencia la necesidad de desarrollar una solución accesible, flexible y 
        personalizable que permita mejorar el control y la trazabilidad ESD en entornos de laboratorio, 
        como el laboratorio TestChip de Intel Costa Rica.

\section{Estructura del Documento}
El documento se organiza de la siguiente manera. En el Capítulo 1 se presenta la introducción, 
los antecedentes y el planteamiento del problema. Posteriormente, 
se desarrolla el marco teórico y el estado del arte asociados a la descarga electrostática, 
la medición de resistencia cuerpo--tierra, la medición de alta impedancia y las soluciones comerciales existentes. 
Luego se describe la solución propuesta, los objetivos del proyecto, el experimento preliminar y el cronograma de trabajo.
